\subsection{先作序关系解读，再作基数解读}

本章记号最安全的用法，是序关系式的：
\begin{itemize}[nosep]
  \item 可以把一个状态看作比另一个状态更具强化性，
  \item 可以把一个目标看作比另一个目标更含糊，
  \item 可以把一种准备机制看作比另一种更充分，
  \item 可以把一次转变看作比另一次更困难。
\end{itemize}

这对本书大多数下游主张来说已经足够。干预设计、案例比较与决策窗口推理，通常更需要排序与变化方向，而不是精确测量。若要作出更强的基数解读，则需要本手稿目前还不具备的校准工作。